\clearpage
\SourceLeaf{486}%
\fancyhead[L]{\small REPORTS ON THE DEGREE SURVEY.}
\pdfbookmark[0]{Report on the results of the trigonometric measurements.}{bericht-resultate-1848}
\begin{center}\large [REPORT ON THE RESULTS OF THE TRIGONOMETRIC MEASUREMENTS.]\end{center}
\par
— — —

Of the calculations used to pass from the raw measurement data [in the journals, which fill 35 notebooks in fair copy] to the final results [in the coordinate register], I have been able to include only a very small portion, in terms of bulk, among the materials I am submitting here.

Indeed, if it were possible to reconstruct all those calculations completely and in full detail, they might easily fill four or six times as many volumes [as the journals now make up in bound form]. But partly because most of the details of those calculations have not been kept at all, and partly because of the form in which the still very bulky bundles of retained papers are found, reconstructing \textit{all} the calculations completely and in order would amount to almost the same thing as doing my entire work over again. I have therefore confined myself to making an orderly extract of the portion of the intermediate calculations that must be regarded as the most telling and useful, namely the tabulation of all the direction angles established and oriented at the various observation stations, with a line parallel to the Göttingen meridian serving as the zero or starting point. These tables have been arranged together in six notebooks under the title \textit{Outlines}\footnote{The total number will be a little over 450.} In this work I received help in several ways from Professor \textsc{Goldschmidt}, who also made most of the fair copies himself. — — —

To account for the very late completion of this work, I must add that the delay arose chiefly because lengthy investigations often had to be made first to resolve the numerous doubts that arose in making the copies, and even more so in preparing the outlines. — — —

As soon as time, health and strength permit, my first two papers on geodetic problems will be followed by a couple of others that go into the particular subjects more closely.

Göttingen, March 15, 1848.
\begin{flushright}C. F. \textsc{Gauss}.\end{flushright}
